\documentclass[11pt]{beamer}
\usepackage{amssymb, latexsym, amsmath, graphics, fullpage, epsfig, amsthm, relsize, pgf, tikz, amsfonts, makeidx, latexsym, ifthen, hyperref, calc}
%\usepackage{eucal} this is a different font than \mathcal{•}
\usetikzlibrary{arrows}
%\newtheorem{theorem}{Theorem}[section]
\newtheorem{remark}{Remark}[section]
%\newtheorem{corollary}{Corollary}[section]
%\newtheorem{lemma}{Lemma}[section]
\newtheorem{assumptions}{Assumptions}[section]
\newtheorem{proposition}{Proposition}
%\newtheorem{definition}{Definition}
\newtheorem{notation}{Notation}
\usepackage{mathrsfs}
%\usepackage[shortlabels]{enumitem}
\usepackage{tikz}
\usepackage{amsmath}
\usetikzlibrary{arrows}
\usetikzlibrary{shadows.blur}
\usepackage{pgfplots}
\usepackage{mwe} % For dummy images
\usepackage{subcaption}
\usepackage{multirow}
\usepackage{dcolumn}
\newcolumntype{2}{D{.}{}{2.0}}
\usepackage{multicol}
\numberwithin{equation}{section}

\newcommand{\lip}{\textup{Lip}_b}
\newcommand{\liph}{\text{Lip}_{\hat\rho}}
\newcommand{\R}{\mathbb{R}}
\newcommand{\E}{\mathbb{E}}
\newcommand{\Norm}[1]{\left\|  #1   \right\|}
\newcommand{\ud}{\ensuremath{\mathrm{d}}}

%\usepackage{lipsum}%% a garbage package you don't need except to create examples.
%\usepackage{fancyhdr}
%\pagestyle{fancy}
%\lhead{\Huge Version A}
%\rhead{\thepage}
%\cfoot{}
%\renewcommand{\headrulewidth}{0pt}



\setbeamersize{text margin left=2mm,text margin right=2mm}

\usepackage{tikz}
\usepackage{pgfplots}
\pgfplotsset{compat=newest}
\usetikzlibrary{decorations.markings}

\allowdisplaybreaks


\defbeamertemplate{footline}{higher page number}
{%
	\hfill%
	\usebeamercolor[fg]{page number in head/foot}%
	\usebeamerfont{page number in head/foot}%
	\insertframenumber\,/\,\inserttotalframenumber\kern1em\vskip20pt%
}
\setbeamertemplate{footline}[higher page number]


%Information to be included in the title page:
\title[About Beamer] %optional
{Invariant Measures for the Stochastic Heat Equation}

\subtitle{Nick Eisenbergs Comprehensive Presentation}

\author[Eisenberg, Chen] % (optional, for multiple authors)
{Nick~Eisenberg \and Le~Chen}

%\institute[VFU] % (optional)
%{
%	\inst{1}%
%	Faculty of Physics\\
%	UNLV
%	\and
%	\inst{2}%
%	Faculty of Chemistry\\
%
%}

\date[Fall 2020] % (optional)
{Fall 2020}

\begin{document}

\frame{\titlepage}

\begin{frame}[t]
\frametitle{Gaussian Processes}

A Gaussian Processes, $G := \{G_t : t \in T\}$ is a type of stochastic processes such that for and finite set of $t_1,...,t_n \in T$, the random vector $(G(t_1),...,G(t_n))$ follows a $n$-dimensional Gaussian distribution. The processes $G$ is uniquely determined by
\begin{enumerate}
		\item $\mu(t) = \E(G(t))$
		\item $C(s,t) = \E(G(t)G(s))$
\end{enumerate}

\begin{theorem}[Important]
	The space of all symmetric and non-negative definite functions, $f(s,t)$, on $T\times T$ matches the space of all covariance functions of all Gaussian processes on $T$. Thus, given a non-negative definite function on $T \times T$ then there exists a Gaussian Processes that possesses this covariance function.
\end{theorem}
\end{frame}

\begin{frame}[t]
\frametitle{Examples of Gaussian Processes}
\begin{enumerate}
	\item \textbf{Brownian Motion:} Consider the following mean and covariance functions on $T=[0, \infty)$.
			\begin{itemize}
				\item $\mu(t) = 0$
				\item $C(s,t) = \min(s,t)$
			\end{itemize}
		It can be shown that $C$ is non-negative definite and the resulting Gaussian Processes is referred to as Brownian Motion, $B= \{ B(t) : t \ge 0 \}$.

	\item \textbf{White Noise:} Consider the following mean and covariance functions on $T=\mathscr{B}(\R^N)$, the Borel sets of $\R^N$.
		\begin{itemize}
			\item $\mu(A) = 0$
			\item $C(A,B) = \lambda^N(A \cap B)$, $\lambda^N$ = $N$-dimensional Lebesgue measure.
		\end{itemize}
	Again it can be shown that $C$ is non-negative definite and thus the resulting Gaussian processes on $T$ is referred to as White Noise, $\dot{W} = \{ \dot{W}(A) : A \in \mathscr{B}(\R^N)  \}$.
\end{enumerate}
\end{frame}

\begin{frame}[t]
	\frametitle{White Noise as an Ito Integral}

First we will briefly explain the meaning of $\int_0^\infty u(t,\omega) \ud B(t,\omega)$, the Ito Integral of $u$, where $B$ is a Brownian Motion.

\begin{definition}[Progressively Measurable]
	A processes, $u(t,\omega)$, is progressively measurable if $u$ restricted to $\Omega \times [0,t]$ is $\mathcal{F}_t \times \mathscr{B}([0,t])$ measurable where $\mathcal{F}_t = \sigma\{B(s) : 0 \le s \le t \}  $.
\end{definition}

\begin{definition}[Progressively Measurable Hilbert space, $L^2(\mathscr{P})$]
We define $L^2(\mathscr{P}) = L^2(\Omega \times [0, \infty), \mathscr{P}, \mathbb{P}\times \lambda ])$ where $\mathscr{P}$ is the sigma-alegebra generated by all the sets $A \subset \Omega \times [0, \infty)$ such that $1_A(t,\omega)$ is progressively measurable. The norm of this Hilbert Space is $\Norm{u}_{\mathscr{P}}^2 = \E \left( \int_0^\infty u(t,\omega)^2 \; \ud t \right)$
\end{definition}

\begin{definition}[Simple Processes and their Ito Integral]
	A process of the form $u(t,\omega) = \sum_{k=1}^n \phi_k(\omega) 1_{(t_k,t_{k+1}]}(t)$ where $\phi_k$ is $\mathcal{F}_k$ measureable is called a simple process. Its Ito integral is defined to be $\int_0^\infty u(t,\omega) \ud B(t,\omega) = \sum_{k=1}^n \phi_k(\omega)(B(t_{k+1},\omega)-B(t_k, \omega))$
\end{definition}
\end{frame}

\begin{frame}[t]
	\frametitle{White Noise as an Ito Integral}
\begin{theorem}
	The simple processes are dense in $L^2(\mathscr{P})$. Thus for any processes, $u(t,\omega)$, there exists a sequence of simple processes, $u_n(t,\omega)$ such that $\E \left( \int_0^\infty (u(t,\omega) - u_n(t,\omega))^2 \ud t \right) \to 0$.
\end{theorem}

\begin{definition}[Ito Integral of a $L^2(\mathscr{P})$ processes]
	For any processes $u(t,\omega) \in L^2(\mathscr{P})$ we define its Ito Integral to be
	\[\int_0^\infty u(t,\omega) \ud B(t,\omega) := \lim_{n \to \infty} \int_0^\infty u_n(t,\omega) \ud B(t, \omega)
	\]
	where the limit is with respect to the $L^2(\mathscr{P})$ norm and the $u_n's$ are from the Theorem.
\end{definition}

\begin{theorem}
	\begin{enumerate}
		\item $ \E \int_0^\infty u(t,\omega) \ud B(t,\omega) = 0$
		\item $ \E\left( \int_0^\infty u(t,\omega) \ud B(t,\omega)\cdot \int_0^\infty v(t,\omega) \ud B(t,\omega) \right) = \E \left( \int_0^\infty u(t,\omega)v(t,\omega) \ud t \right) $
	\end{enumerate}
\end{theorem}
\end{frame}

\begin{frame}[t]
	\frametitle{White Noise as an Ito Integral}
With the above in consideration, we define a set function on $\mathscr{B}(\R)$
\[
	\dot{W}(A) = \int_0^t 1_A(t) \ud B(t,\omega)
\]
For any borel set, $A$, the deterministic indicator function $1_A \in L^2(\mathscr{P})$ and so the Ito Integral exists.

From the theorem in the previous slide we see that
\begin{itemize}
	\item $\E\dot{W}(A) = 0 $
	\item $ \E(\dot{W}(A) \dot{W}(B) = \int_0^\infty 1_A(t)1_B(t) \ud t = \lambda (A \cap B) $
\end{itemize}
Thus, $\dot{W}$ is a white noise on $\mathscr{B}(\R)$

\end{frame}

\begin{frame}[t]
	\frametitle{Walsh Integral}
The Walsh integral is a generalization of the Ito integral and is used to study the stochastic heat equation with either a space-time white noise force term in the one dimensional case or a colored noise in higher dimensions which we will define below. The stochastic heat and wave equation with white noise forcing term only exists in dimension 1. For that we reason we define the colored noise which is a little more smooth.

The Walsh Integral is constructed in a similar fashion as the Ito integral. Where as the integrator in the Ito integral is a Brownian Motion, the integrator in the Walsh integral is a worthy martingale measure. It wont be really necessary to discuss the definition of the integrator. But if we consider a white noise on $\mathscr{B}(\R^N)$, $\dot{W}$, then for $A \in \mathscr{B}(\R^{N-1})$
\[
	W_t(A) = \dot{W}([0,t] \times A)
\]
is a worthy martingale measure and can be used as an integrator for the Walsh integral.
\end{frame}

\begin{frame}[t]
	\frametitle{Walsh Integral}
As in the Ito case we consider simple functions which are linear combinations of functions of the form $f(t,x,\omega) = X(\omega)1_{(a,b]}(t)1_A(x)$. We then let $\mathscr{P}$ denote the sigma-field generated by all the simple processes. Then given a worthy martingale measure, we can define the walsh integral
\[
\int_0^t \int_A f(s,y) M(\ud s, \ud y)
\]
 for any $\mathscr{P}$ measurable function satisfying for all $T>0$
 \[
 	\int_0^T\iint_{\R^{2n}}|f(x,t)f(y,t)|K(\ud x, \ud y, \ud t) < \infty
 \]
where $K$ is defined through the worth martingale measure. We define this integral through a limit of Walsh integrals of simple processes as in the Ito case.

\end{frame}

\begin{frame}[t]
	\frametitle{Colored Noise}
	We first start with a non-negative and non-negative definite tempered measure $\Gamma$ with density $f$. In other words, for any $\varphi \in \mathscr{S}(\R^d)$ we have
	\[
		\int_{\R^d} \Gamma(\ud x) (\varphi * \tilde{\varphi})(x) \ge 0, \quad \tilde{\varphi}(x) = \varphi(-x)
	\]
 Bochner's Theorem says that there exists a positive measure $\mu$ such that $\int \Gamma(\ud x) \varphi(x) = \int \mu(\ud x) \mathcal{F}\varphi(x)$ and $\mu$ is referred to as the spectral measure and plays a key role in discussing the solveability of an SPDE. We define on 	$C_0^{\infty}(\R^{d+1})$ the functional
	\begin{align*}
	C(\varphi, \psi) &= \int_0^\infty \ud s \int_{\R^d} \Gamma(\ud x) (\varphi(s,\cdot)*\tilde{\psi}(s,\cdot))(x)
	\\& = \int_0^\infty \ud s \int_{\R^d} \ud x \int_{\R^d} \ud y \varphi(s,x)f(x-y)\psi(s,y)
	\end{align*}
 which is positive definite. Thus we can assosciate this with the covariance of a $0$-mean Gaussian Processes, $F= \{ F(\varphi) : \varphi \in C_0^\infty (\R^{d+1})\}$
\end{frame}

\begin{frame}[t]
	\frametitle{Walsh integral with Colored Noise Integrator}
	We need to construct the colored noise processes to use as an integrator. Let $A \in \mathscr{B}(\R^d)$ and take functions $\varphi_n \in \mathscr{S}(\R^d)$ such that $\varphi_n \to 1_A$. We can in fact show that $\E[(F(\varphi_n) - F(\varphi_m))^2] \to 0$ and we denote the $L^2(\Omega)$ limit as $F(A)$. Thus we may define
	\[
		M_t(A) = F([0,t]\times A)
	\]
	It turns out that $M$ is a worthy martingale measure and through construction of the Walsh integral we deduce that
	\[
		F(\varphi) = \int_0^\infty \int_{\R^d} \varphi(s,x) M(\ud x, \ud s)
	\]
	It can be shown that the suitable integrands must satisfy
	\[
		\int_0^\infty \ud s \int_{\R^d} \ud x \int_{\R^d} \ud y \varphi(s,x) f(x-y) \varphi(s,y) < \infty
	\]
\end{frame}

\begin{frame}[t]
	\frametitle{Stochastic Heat Equation}
	\[
	 \begin{cases}
	\dfrac{\partial u}{\partial t}(t,x)- \frac{1}{2}\Delta_xu(t,x) = \sigma(u(t,x)\dot{M}(t,x) &
	\\ u(0,\cdot) = \mu(\cdot)  &
	\end{cases}
	\]
Above is the informal representation of the SHE. We say informal because $\dot{M}(t,x)$ does not exist. $\dot{M}(t,x)$ is acting as the density of $M(\ud s, \ud x)$ in the sense that
\[
	\int_0^t \int_{\R^d}\varphi(s,y) M(\ud s, \ud y) = \int_0^t \int_{\R^d}\varphi(s,y)\dot{M}(s, y)\ud s \ud y
\]
but this does not exist. The rigorous meaning of the SHE is the following integral equation
\[
	u(t,x) = \int_{\R^d} G(t,x-y) \mu(\ud y) + \int_0^t \int_{\R^d}G(t-s,x-y)\sigma(u(s,y))M(\ud s, \ud y)
\]
where $G(t,x)$ is the heat kernel. Under certain conditions on the spectral measure, initial condition and $\sigma$, there exists a unique solution
\end{frame}

\begin{frame}[t]
	\frametitle{Definition of Invariant Measures}
	Consider an arbitrary measure space $(E,\mathscr E)$ and denote $\mathcal M_1(E)$ to be the space of all probability measure on $(E,\mathscr E)$. Before we give the definition of an invariant measure, we need to define what a Markovian Transition Function is.
	\begin{definition}[Markovian Transition Function]  We say that $P_t(x,\Gamma)$, $t>0$, $x\in E$, $\Gamma \in \mathcal E$, is a Markovian Transition Function is:

		\begin{enumerate}[(i)]
			\item $P_t(x,\cdot)$ is a probability measure on $(E,\mathscr E)$ for $t \ge 0$, $x \in  E$
			\item $P_t(\cdot,\Gamma)$ is a $\mathscr E$-measurable function for $t \ge 0$ and $\Gamma \in \mathscr E$
			\item $P_{t+s}(x,\Gamma) = \int_EP_t(x,\ud y)P_s(y,\Gamma)$
			\item $P_0(x,\Gamma)=1_\Gamma(x)$ where $1_\Gamma(x)$ is the indicator function.
		\end{enumerate}
	\end{definition}
	Informally, $P_t(x,\Gamma)$ can be thought of as the probability that $x$ moves into $\Gamma$ after $t$ amount of time.
\end{frame}

\begin{frame}[t]
	\frametitle{Definition of Invariant Measures}
	\begin{definition} [Invariant Measure]
		We say that a measure $\eta(\cdot) \in \mathcal M_1(E)$ is invariant for transition functions $\{P_t(x,\Gamma)\} _{t\ge 0}$ if for any $\Gamma \in \mathscr E$ we have that $$ \eta(\Gamma) = \int_E P_t(x,\Gamma)\eta(\ud x) \quad \text{for any $t\ge 0$} $$
	\end{definition}

	Recall from Calculus that the average of a function $f:[a,b] \to \R$ is defined as
	$$ \text{AVG}(f) = \dfrac{1}{a-b}\int_a^b f(x) \ud x$$
	From this we see that an invariant measure is just the average of the transition functions. Namely, if $\eta$ is invariant for $P_t(x,\Gamma)$ then $\eta(\Gamma)$ is the average of the function $ x \mapsto P_t(x,\Gamma)$ when considering the measure space $(E,\mathscr E, \eta)$.

\end{frame}

\begin{frame}[t]
	\frametitle{Main Result}
	The laws of the solution the the SHE defines Markovian transition functions
	\[
	P_t(\varphi , \Gamma) = \mathscr L (u^\varphi(t,\cdot))(\Gamma) = P[\omega \in \Omega : u^\varphi(t,\cdot)\in \Gamma] \quad \Gamma \in \mathscr B (L^2_\rho)
	\]
	where $u^\varphi(t,x)$ is the solution to the SHE starting at $\varphi \in L^2_\rho$. We have the following main results
	\begin{theorem}
		Under certain conditions of the spectral density $\hat{f}$ and on the admissible weight $\rho(x)$, there exists a measure $\eta$ that is invariant for the solution to the stochastic heat equation for any $L^2_\rho$ initial condition. In other words, for $P_t(\varphi,\Gamma)$ we have that, $ \eta(\Gamma) = \int_{L^2_\rho(\R^d)}P_t(x,\Gamma)\eta( \ud x) $ for $ t\ge T >0$ for any $T>0$.
	\end{theorem}

	\begin{theorem}
		Under the same assumptions as the theorem above, there exists a measure $\eta$ that is invariant for the solution to the stochastic heat equation with the $\delta_0(x)$ initial condition.
	\end{theorem}

\end{frame}

\begin{frame}[t]
	\frametitle{Improvements to Previous Results}
	This result on the existence of invariant measures serves two purposes.

	It first acts as a translation of the paper \cite{TessitoreZabczyk98M} from Daprato theory to Walsh Theory. There are two main approaches to defining the stochastic integral, the first being the theory laid out by Daprato and the second by Walsh. Both theories have been shown to be equivalent, however some problems may be handled more easily using one theory over the other. In fact, you can see from the proof of Lemma 3.3 in mine and Dr. Chen's paper , which mimics Theorem 3.1 in \cite{TessitoreZabczyk98M}, that proving the result using the Theory of Walsh requires much more work. For that reason, the vast majority of all research done on invariant measures happens to be done using the theory of Daprato.
\end{frame}

\begin{frame}[t]
	\frametitle{Improvements to Previous Results}
	Because of using the theory of Walsh, we are able to apply the estimates on the moments of the solution to the stochastic heat equation proven in \cite{CH16Comparison}. These moment estimates allowed us to improve previous results in the following ways:
	\begin{enumerate}
		\item We may show the existence of invariant measures for the SHE with measure valued initial conditions such as the dirac delta distribution.
		\item We may show the existence of invariant measures for the SHE with an $L_\rho^2(\R^d)$ function where as in previous results, such as in \cite{MSY16, TessitoreZabczyk98M, AM2002}, the initial condition either had to be bounded and continuous or belong to a specific $L_\rho^2(\R^d)$. In \cite{AM2002 } they allow for a non-zero and unbounded drift term which, at the moment, can not be handled by the theory of Walsh. So this further shows the trade off between the two theories of the stochastic integral.
	\end{enumerate}
\end{frame}

\begin{frame}[t]
	\frametitle{Current Project}
	The new project Dr. Chen and I are working on deals with the following fractional SPDE
	\begin{equation}
	\begin{cases}
		\left(\partial^b + \frac{\nu}{2}(-\Delta)^{a/2}\right)u(t,x) = I^r_t \left[\rho(u(t,x))\dot W(x) \right] & \text{$x\in \R^d$ , $t>0$} \\
		u(0,\cdot) = 1                                                                                           & \beta \in (0,1]            \\
		u(0,\cdot) = 1, \quad u_t(0,\cdot) = 0                                                                   & \beta \in (1,2]
	\end{cases}  \label{fspde}
	\end{equation}
	where $a \in (0,2]$, $b \in (0,2)$, $\nu>0$, $r>0$ and $\theta >0$. In the equation, $(-\Delta)^{a/2}$ refers to the fractional Laplacian \cite{FPDE} and $\partial^b$ refers to the Caputo derivative \cite{PodFDE} and $I^r_t$ is the Riemann-Liouville fractional integral operator.

	The noise $W:= \{ W(\varphi) : \varphi \in \mathcal{D}(\R^d) \}$ in this equation is a time-independent Gaussian iso-normal processes defined on the  space of test-functions with covariance
	\[
		\mathbb{E}(W(\varphi)W(\psi)) = \int_{\R^d}\mathcal{F}\varphi(\xi)\overline{\mathcal{F}\psi}(\xi) \mu(\ud \xi)
	\]
	where $\mu$ is the spectral measure from Bochner's Theorem.
\end{frame}

\begin{frame}[t]
\frametitle{Current Project}
The solution of the fractional equation is defined through the Skorohod Integral \cite{NauMal}, which is a generalization of the Ito Integral. Because of the constant one initial conditions, we say that $u$ solves the fractional equation if $u$ satisfies the following integral equation
\[
	u(t,x) = 1 + \sqrt{\theta}\int_0^t \left( \int_{\R^d} G(t-s,x-y) W(\ud y)  \right) \ud s
\]
Through a standard Piccard-iteration scheme, we can show that the solution has a chaos-expansion
\[
u(t,x) = 1 + \sum_{n \ge 1} \theta^{n/2}I_n(f_n(\cdot,x,t))
\]
where $I_n$ is the $n$ dimensional Ito Integral with respect to $W$ and the kernals $f_n$ are given by
\[
f_n(x_1,...,x_n,x,t) = \int_{[0,t]^n<} G(t-t_n,x-x_n)\cdots G(t_2-t_1,x_2-x_1) \ud \textbf{t}
\]
\end{frame}

\begin{frame}[t]
	\frametitle{Current Project}
	Things to note about the fractional SPDE is that when $b=1$ and $r=0$ and $a=2$ then we recover the SHE. When  $b=2$ and $r=0$ and $a=2$ then we recover the SWE.

	The goal of this project is to prove the existence and uniqueness of the solution to the fractional equation as well as prove some asymptotics which have already been proven for the SWE in \cite{LeXiaReb}.

	We have already shown existence and uniqueness to the fractional equation, which followed by mimicking the proof in \cite{LeXiaReb}.

\end{frame}

\begin{frame}[allowframebreaks]
	\frametitle<presentation>{References}
	\begin{thebibliography}{10}

	\bibitem{AM2002}
	Assing, Sigurd and Manthey, Ralf.
	\newblock Invariant measures for stochastic heat equations with unbounded coefficients
	\newblock {\em Stoc. Proc. and their App.}  2002.

	\bibitem{LeXiaReb}
	Balan, R., Chen, Le and Chen, Xia.
	\newblock Exact asymptotics of the stochastic wave equation with time-independent noise.

	\bibitem{CH16Comparison}
	Chen, Le and Jingyu Huang.
	\newblock Comparison principle for stochastic heat equation on $\R^d$.
	\newblock {\em Ann.\ Probab.}  2019, to appear.

	\bibitem{CK15SHE}
	Chen, Le and Kunwoo Kim.
	\newblock Nonlinear stochastic heat equation driven by spatially colored noise: moments and intermittency.
	\newblock {\em Acta Math. Sci. Ser. B}, 2019, to appear.

	\bibitem{DalangEtc09Minicourse} Dalang, Robert, Davar Khoshnevisan, Carl
	Mueller, David Nualart, and Yimin Xiao
	\newblock {\it A minicourse on stochastic partial differential equations}.

	\bibitem{DZ92}
	Da Prato, Giuseppe and Jerzy Zabczyk.
	\newblock {\it Stochastic equations in infinite dimensions.}
	Encyclopedia of Mathematics and its Applications, 44. Cambridge University Press, Cambridge, 1992. xviii+454 pp.


	\bibitem{DZ96}
	Da Prato, Giuseppe and Jerzy Zabczyk.
	\newblock {\it Ergodicity for infinite-dimensional systems. }
	\newblock London Mathematical Society Lecture Note Series, 229. Cambridge University Press, Cambridge, 1996.

	\bibitem{FPDE}
	Folland, Gerald
	\newblock Introdution to Partial Differential Equations.

	\bibitem{MSY16}
	Misiats, Oleksandr, Oleksandr Stanzhytskyi, and Nung-Kwan Yip.
	\newblock Existence and uniqueness of invariant measures for stochastic reaction-diffusion equations in unbounded domains.
	\newblock {\it J. Theoret. Probab.} 29 (2016), no. 3, 996--1026.

	\bibitem{NauMal}
	Nualart, D and Nualart, E.
	\newblock Introduction to Malliavin Calculus
	\newblock {\it Cambridge. Cambridge University Press} (2018)

	\bibitem{PodFDE}
	Podlubny, L.
	\newblock Fractional Differential Equations
	\newblock {\it Mathematics in Science and Engineering}


	\bibitem{TessitoreZabczyk98M}
	Tessitore, Gianmario and Jerzy Zabczyk.
	\newblock Invariant measures for stochastic heat equations.
	\newblock {\it Probab. Math. Statist.} 18 (1998), no. 2, Acta Univ. Wratislav. No. 2111, 271--287.

	\bibitem{Walsh} Walsh, John B.
	\newblock {\it An Introduction to Stochastic Partial Differential Equations}.
	\newblock In: \`Ecole d'\`et\'e de probabilit\'es de
	Saint-Flour, XIV---1984, 265--439.
	Lecture Notes in Math.\ 1180, Springer, Berlin, 1986.
	\end{thebibliography}
\end{frame}



\end{document}
